% -*- mode: LaTeX; TeX-PDF-mode: t; -*-
\input{@resources/tex-add-search-paths}  % allow latex to find custom stuff
\input{./.econtexRoot}\documentclass[\econtexRoot/BufferStockTheory]{subfiles}

\newcommand{\subname}{ApndxHarKmenberg}
%\input{\econtexRoot/@local/dir-paths} % get directory macros
\usepackage{econark-ifsubfile}        % allow conditional execution of code
\usepackage{econark-xrsetup}          % Xternal crossReferences (from main document)

\providecommand{\koppa}{k}
\xrsetup{\econtexRoot/\texname}

  %\externaldocument{BufferStockTheory}\providecommand{\texname}{}\renewcommand{\texname}{Introduction}} % Get xrefs -- esp to appendix -- from main file; only works properly if main file has already been compiled;
\begin{document}

\begin{comment}
Consider any configuration of parameter values under which, with the variance of permanent shocks equal to zero, the {\GICRaw} holds exactly:
\begin{equation}\begin{gathered}\begin{aligned}
  \GICRaw & = 1 = \exp\left(\tilde{\Ex}[\log \permShk]\right)
\end{aligned}\end{gathered}\end{equation}
\end{comment}

\section{Harmenberg's Method}\label{sec:ApndxHarKmenberg}

\newcommand{\permLvlPrb}{p}
\newcommand{\permLvlVec}{\mathrm{p}}
\newcommand{\lomdkPrb}{\pi} 
\newcommand{\lomdkPrbmm}{\varpi} 
\newcommand{\lomdkMat}{\Pi}
\newcommand{\Prb}{f}
%\newcommand{\Vct}{\mathrm{F}}
\newcommand{\Vct}{\digamma}
\newcommand{\permShkPrb}{\Prb_{\permShk}}
\newcommand{\permShkVec}{\Vct_{\permShk}}
\newcommand{\permShkRawPrb}{\Prb_{\permShk}}
\newcommand{\permShkRawVec}{\Vct_{\permShk}}
\newcommand{\permShkWgtPrb}{\tilde{\Prb}_{\permShk}}
\newcommand{\permShkWgtVec}{\tilde{\Vct}_{\permShk}}
\newcommand{\mpPrb}{\chi}
\newcommand{\mpPrbMarg}{\mpPrb^{\mNrm}}
\newcommand{\Wgt}{\tilde}
\newcommand{\Raw}{}
\newcommand{\mpPrbMargRaw}{\Raw{\mpPrb}^{\mNrm}}
\newcommand{\mpPrbMargWgt}{\Wgt{\mpPrb}^{\mNrm}}
\newcommand{\mpMat}{\mathrm{X}}
\newcommand{\mpMatMarg}{\mpMat^{\mNrm}}
\newcommand{\mpMatMargRaw}{\Raw{\mpMat}^{\mNrm}}
\newcommand{\mpMatMargWgt}{\Wgt{\mpMat}^{\mNrm}}
\newcommand{\mNrmVec}{\mathrm{m}}
\newcommand{\mNrmNow}{\mNrm_{t}}
\newcommand{\mNrmNxt}{\mNrm_{t+1}}
\newcommand{\Bot}{\underline}
\newcommand{\mNrmNxtBot}{\Bot{\mNrm}_{t+1}}
\newcommand{\permShkNow}{\permShk_{t}}
\newcommand{\permShkPrbNxt}{\permShk_{t+1}}
\newcommand{\permShkPrbBot}{\Bot{\permShk}}
\newcommand{\permShkPrbNxtBot}{\Bot{\permShk}_{t+1}}
\newcommand{\permLvlPrbNxt}{\permLvlPrb_{t+1}}
\newcommand{\permLvlPrbNxtBot}{\Bot{\permLvlPrb}_{t+1}}
\newcommand{\permLvlPrbNow}{\permLvlPrb_{t}}
\newcommand{\tranShkNow}{\tranShk_{t}}
\newcommand{\tranShkNxt}{\tranShk_{t+1}}
\newcommand{\mLvlNow}{\mLvl_{t}}
\newcommand{\mLvlNxt}{\mLvl_{t+1}}

Harmenberg defines a `density kernel' describing the law of motion for the normalized state $\lomdkPrb(\mNrmNxt,\mNrmNow,\permShkPrbNxt)$ and defines $\permShkRawPrb$ as the density of the permanent shock distribution; we correspondingly define $F_{\permShk}$ as the CDF of permanent shocks.

The joint cumulative distribution boldface $\pmb{\mpPrb}$ for $\mNrmNxt$ and $\permLvlPrbNxt$ as a function of the stochastic variables and the joint marginal distributions (nonbold $\mpPrb$) in period $t$ is:
\begin{equation}
  \begin{aligned} \label{eq:MeasTrnstCDF}
  \pmb{\mpPrb}_{t+1}(\mNrmNxt,~& \permLvlPrbNxt) =  \\ &
                                          \int_{\mNrmNxtBot}^{\mNrmNxt}
                                          \int_{\permLvlPrbNxtBot}^{\permLvlPrbNxt}
                                          \int_{\permLvlPrbNow} 
                                          \int_{\mNrmNow} \lomdkPrb_{t}\left(\mNrmNxt,\mNrmNow,\permShkPrbNxt\right)\mpPrb_{t}(\mNrmNow,\permLvlPrbNow)  d \mNrmNow                                                                                     d \permLvlPrbNow
                                                d F_{\permShk(\permLvlPrbNxt)}
                                                d \mNrmNxt \notag
\end{aligned}\end{equation}
where 
\begin{align*}
  \permShkPrbNxt & = \permLvlPrbNxt/(\PermGroFac \permLvlPrbNow) 
\\    F_{\permShk(\permLvlPrbNxt)} & = \int_{\permShkPrbBot}^{\permShk(\permLvlPrbNxt)} f_{\permShk} d\permShk
    \\ dF_{\permShk(\permLvlPrbNxt)} & = \left(\frac{d (\permLvlPrbNxt/(\PermGroFac \permLvlPrbNow))}{d\permLvlPrbNxt}
                             \right)f_{\permShk}(\permShkPrbNxt) d \permLvlPrbNxt 
\\ & = \left(\frac{1}{\PermGroFac \permLvlPrbNow}                             \right)f_{\permShk}(\permShkPrbNxt) d \permLvlPrbNxt 
  \end{align*}

\cite{harmenbergInvariant}'s first equation corresponds to the marginal density obtained by differentiating the above CDF with respect to $\mNrmNxt$ and $\permLvlPrbNxt$ and says that the `Markov operator that maps a distribution $\mpPrb \in D(\mNrm \times \permLvlPrb)$ to the next-period $\mpPrb$' (that is, defines the dynamics of the joint marginal distribution of $\mNrm$ and $\permLvlPrb$) is given by
\begin{equation}\begin{gathered}\begin{aligned} \label{eq:MeasTrnst}
\mpPrb_{t+1}(\mNrmNxt,\permLvlPrbNxt) & =
                                          \int_{\permLvlPrbNow} \left[
                                           \int_{\mNrmNow} \lomdkPrb_{t}\left(\mNrmNxt,\mNrmNow,
                                           \overbrace{\frac{\permLvlPrbNxt}{\PermGroFac \permLvlPrbNow}}^{\permShkPrbNxt}\right)\permShkRawPrb(\frac{\permLvlPrbNxt}{\PermGroFac \permLvlPrbNow})\mpPrb_{t}(\mNrmNow,\permLvlPrbNow)\left(\frac{1}{\PermGroFac \permLvlPrbNow}\right)  d \mNrmNow
                                          \right] d \permLvlPrbNow
\end{aligned}\end{gathered}\end{equation}
\newcommand{\nPermShk}{\texttt{i}}\newcommand{\nmNrmVecNow}{\texttt{j}}\newcommand{\nmNrmVecNxt}{\texttt{k}}\newcommand{\permLvlNow}{\texttt{n}}\newcommand{\npermLvlNxt}{\texttt{q}}
\newcommand{\npermLvlNow}{\digamma} % to bring attention to it as something that needs to be fixed in next revision
\newcommand{\npermLvlVecNow}{\texttt{n}}\newcommand{\npermLvlVecNxt}{\texttt{q}}\newcommand{\nPermShkNxtFunc}{\iota(\npermLvlVecNow,\npermLvlVecNxt)}\newcommand{\nPermShkVecNxt}{\texttt{i}}
\newcommand{\mNrmVecNow}{\mNrmVec_{t}^{\nmNrmVecNow}}
\newcommand{\mNrmVecNxt}{\mNrmVec_{t+1}^{\nmNrmVecNxt}}
\newcommand{\permShkVecNxt}{\permShk_{t+1}^{\nPermShkVecNxt}}
\newcommand{\permShkVecNow}{\permShk_{t}^{\nPermShkVecNow}}
\renewcommand{\permLvlNow}{\permLvlPrb_{t}}
\newcommand{\permLvlNxt}{\permLvlPrb_{t+1}}
\newcommand{\permLvlVecNow}{\permLvlVec_{t}^{\npermLvlVecNow}}
\newcommand{\permLvlVecNxt}{\permLvlVec_{t+1}^{\npermLvlVecNxt}}

\noindent A somewhat awkward notational scheme allows us to define an almost completely parallel representation of the corresponding discrete transition process:
\begin{enumerate}
\item $\{\nPermShk, \nmNrmVecNow, \nmNrmVecNxt\}$ index elements of vectors identifying possible values of $\permShkPrbNxt$, $\mNrmNow$, and $\mNrmNxt$
  \begin{itemize}
  \item We use the capital of the Roman letter to count the number of possible entries
    \item e.g., there are $K$ possible different values for $\mNrmVec:\{\mNrmVec_{t}^{0},\mNrmVec_{t}^{1},...,\mNrmVec_{t}^{K}\}$
    \end{itemize}
\item $\npermLvlVecNow$ indexes the level of permanent income now: $\permLvlVec_{t}[\npermLvlVecNow]$
\item$\lomdkMat_{t}[\nmNrmVecNxt,\nmNrmVecNow,\nPermShk]$ indicates the probability of making a transition from value $\nmNrmVecNow$ of $\mNrmNow$ to value $\nmNrmVecNxt$ of $\mNrmNxt$ given that the realization of $\permShkPrbNxt$ is $\permShkPrbNxt[\nPermShk]$
\item $\permShkVec[\nPermShk]$ is the probability of drawing the $\nPermShk$'th value of $\permShkPrbNxt$
\end{enumerate}
then for a person whose location is described in period $t$ as being at permanent income level $\permLvlVecNow$ and market resources ratio $\mNrmNow[\nmNrmVecNow]$, the elements of the matrix in the next period are given by:
\begin{equation}\begin{aligned} \label{eq:mpMattp1}
  \mpMat_{t+1}[\nmNrmVecNxt,\npermLvlVecNxt] & = \sum_{\npermLvlVecNow} \left(\sum_{\nmNrmVecNow} \lomdkMat_{t}[\nmNrmVecNxt,\nmNrmVecNow,\nPermShkNxtFunc]\permShkVec[\nPermShkNxtFunc]\mpMat_{t}[\nmNrmVecNow,\npermLvlVecNow]  \right)
\end{aligned}\end{equation}
where $\mNrmVec_{t+1}[\nmNrmVecNxt]$ is the $\nmNrmVecNxt$'th element of the vector $\mNrmVec_{t+1}$, $\permLvlVec_{t+1}[\npermLvlVecNxt]$ is the $\npermLvlVecNxt$'th element of $\permLvlVec_{t+1}$, and $\nPermShkNxtFunc$ is a function that calculates the index value of $\nPermShk$ that would achieve the transition from $\permLvlVecNow$ to $\permLvlVec_{t+1}[\npermLvlVecNxt]$.\footnote{Of course, unless the is highly unlikely that $\nPermShkNxtFunc$ would yield an integer unless the elements of $\npermLvlVecNxt$ were chosen as the unique elements defined by all possible combinations of $\permLvlVec_{t}$ and $\permShkVec$; one option is to allocate probabilities in proportion to distance to the nearest integer values of $\nPermShkNxtFunc$ above and below the current point.
(There are other options, which may not be meaningfully better).}$^{,}$\footnote{The $1/\PermGroFac\permLvlNow$ is present in the continuous version but not the discrete version because the latter captures masses and the former captures densities.}
Harmenberg defines 
\begin{equation}\begin{gathered}\begin{aligned}
     \mpPrbMargRaw_{t}(\mNrmNow) := & \int_{\permLvlPrbNow} \mpPrb_{t}(\mNrmNow,\permLvlPrbNow) d \permLvlPrbNow
\end{aligned}\end{gathered}\end{equation}
which measures the population density of persons whose market resources are $\mNrmNow$.
In matrix terms, the corresponding representation is:
\begin{equation}\begin{gathered}\begin{aligned}
     \mpMatMargRaw_{t}[\nmNrmVecNow] = & \sum_{\npermLvlVecNow} \mpMat_{t}[\nmNrmVecNow,\npermLvlVecNow] 
\end{aligned}\end{gathered}\end{equation}
which makes it easy to see that $\mpMatMargRaw_{t}[\nmNrmVecNow]$ just measures the total probability mass associated with all possible levels of permanent income for agents at $\mNrmVec_{t}^{\nmNrmVecNow}$.
That is, it tells us \textit{how many agents} have $\mNrmNow=\mNrmVecNow$.

In order to compute the absolute aggregate \textit{amount} of market resources $\mLvlNow$ (boldface indicates levels) owned by people with a market resources \textit{ratio} of $\mNrmVecNow$, we need to know the total \textit{amount} of permanent income accruing to those people:
\begin{equation}\begin{gathered}\begin{aligned}
  \mLvlNow[\nmNrmVecNow] & = \mNrmVecNow \underbrace{\sum_{\npermLvlVecNow} \permLvlVecNow\mpMat_{t}[\nmNrmVecNow,\npermLvlVecNow]}_{\equiv  \mpMatMargWgt[\nmNrmVecNow]\PermGroFac^{t}}
\end{aligned}\end{gathered}\end{equation}
where $\mpMatMargWgt[\nmNrmVecNow]$ is what Harmenberg calls the Permanent Income Weighted measure.
Under the assumption that aggregate permanent income was 1.0 in period 0 and has grown by the factor $\PermGroFac$ thereafter, the following direct formula for $\mpMatMargWgt$ can be seen to capture the \textit{proportion of aggregate permanent income earned} by people at the given $\mNrmNow$:
\begin{equation}\begin{gathered}\begin{aligned}
  \mpMatMargWgt[\nmNrmVecNow] & = \PermGroFac^{-t} \sum_{\npermLvlVecNow} \permLvlVecNow\mpMat_{t}[\nmNrmVecNow,\npermLvlVecNow]
\end{aligned}\end{gathered}\end{equation}



With this in hand, it is a simple matter to compute the total aggregate mass of $\MLvl$:
\begin{equation}\begin{gathered}\begin{aligned}
  \MLvl_{t} & = \PermGroFac^{t} \sum_{\nmNrmVecNow} \mNrmVecNow \mpMatMargWgt[\nmNrmVecNow]
\end{aligned}\end{gathered}\end{equation}
or for that matter consumption:
\begin{equation}\begin{gathered}\begin{aligned}
  \CLvl_{t} & = \PermGroFac^{t} \sum_{\nmNrmVecNow} \cFunc(\mNrmVecNow) \mpMatMargWgt[\nmNrmVecNow]
\end{aligned}\end{gathered}\end{equation}

Harmenberg formulates his corresponding propositions in the continuous description of the problem using probability measures, e.g.:
\begin{equation}\begin{gathered}\begin{aligned}
  \mpPrbMargWgt_{t}(\mNrmNow) & = \int_{\permLvlNow} \permLvlNow \mpPrb(\mNrmNow,\permLvlNow) d \permLvlNow
\\ \CLvl_{t} & = \PermGroFac^{t} \int_{\mNrmNow} \cFunc(\mNrmNow) \mpPrbMargWgt(\mNrmNow)                             d \mNrmNow
\end{aligned}\end{gathered}\end{equation}

The point here is `that the permanent-income-weighted distribution is a sufficient statistic for computing aggregate consumption, aggregate savings, and similar aggregate variables'.
Thus, rather than requiring us to keep track of the multidimensional joint distribution over $\mNrm$ and $\permLvl$, we need only know the distribution of permanent-income-weighted $\mNrm$.

The crucial last step is to define the law of motion for the weighted system.
In the continuous formulation, Harmenberg shows (his Theorem 1) that, if we define a `permanent-income-shock-weighted' version of the original permanent shock distribution\footnote{(Harmenberg calls this the `permanent-income-neutral measure,' which is slightly confusing as it does not involve the level of permanent income but only the shocks thereto).} as
\begin{equation}\begin{gathered}\begin{aligned}
  \permShkWgtPrb(\permShkPrbNxt) & = \permShkPrbNxt \permShkRawPrb(\permShkPrbNxt)
\end{aligned}\end{gathered}\end{equation}
then the laws of motion are respectively given by 
\begin{equation}\begin{aligned}
  \mpPrbMargRaw_{t+1}(\mNrmNxt) & = \int_{\mNrmNow} \int_{\permShkPrbNxt} \lomdkPrb(\mNrmNxt,\mNrmNow,\permShkPrbNxt) \mpPrbMargRaw_{t}(\mNrmNow) \permShkRawPrb(\permShkPrbNxt) d \permShkPrbNxt d \mNrmNow
  \\
  \mpPrbMargWgt_{t+1}(\mNrmNxt) & = \int_{\mNrmNow} \int_{\permShkPrbNxt} \lomdkPrb(\mNrmNxt,\mNrmNow,\permShkPrbNxt) \mpPrbMargWgt_{t}(\mNrmNow) \permShkWgtPrb(\permShkPrbNxt) d \permShkPrbNxt d \mNrmNow
\end{aligned}\end{equation}
(where the difference between the two is the presence or absence of the $\sim$ accent in three places).

The key steps in the proof are the change in variables in which $\permShk_{t+1}=\permLvlNxt/(\PermGroFac \permLvlNow)$ and a change in the order of integration which is permitted by \href{https://en.wikipedia.org/wiki/Fubini\%27s_theorem}{Fubini's theorem}.

Omitting the $\PermGroFac$ term (equivalently, setting it to $\PermGroFac=1$), the discrete version of the proof is \renewcommand{\PermGroFac}{}
\begin{equation}\begin{gathered}\begin{aligned}
  \mpMatMargWgt[\nmNrmVecNxt] & = \sum_{\npermLvlVecNxt} \permLvlVecNxt\mpMat_{t+1}[\nmNrmVecNxt,\npermLvlVecNxt]
\\   & \approx \sum_{\npermLvlVecNxt} \sum_{\nmNrmVecNow} \sum_{\npermLvlVecNow} \permLvlVecNxt \lomdkMat(\nmNrmVecNxt,\nmNrmVecNow,\npermLvlVecNxt) \permShkWgtVec(\nPermShkNxtFunc)  \mpMat_{t}(\nmNrmVecNow,\npermLvlVecNow)
%  \\   & = \sum_{\nmNrmVecNow} \sum_{\npermLvlVecNxt} \sum_{\npermLvlVecNow} \permLvlVecNxt \lomdkMat(\nmNrmVecNxt,\nmNrmVecNow,\npermLvlVecNxt) \permShkWgtVec(\nPermShkNxtFunc) \mpMat_{t}(\nmNrmVecNow,\npermLvlVecNow)/\permLvlVecNow
\end{aligned}\end{gathered}\end{equation}
% where the last step corresponds to the change in the order of integration in Harmenberg's proof (
and the $\approx$ captures the fact that in the discrete context the necessity to allocate masses to points on the grid will lead to approximation error.

The change of variables is accomplished by realizing that just as there was an $\iota$ that gave us the appropriate $\permShk$ required to get from $\permLvlNow$ to $\permLvlNxt$, we can define a $\koppa(\npermLvlNow,\nPermShkVecNxt)$ that lets us approximate the $\npermLvlNxt$ needed as an argument to $\lomdkMat$ and the index for $\permLvlVecNxt$.
Now we can sum over the permanent shocks indexed by $\nPermShkVecNxt$:
\begin{equation}\begin{gathered}\begin{aligned}
  \mpMatMargRaw[\nmNrmVecNxt] & =
                                \sum_{\nmNrmVecNow}
                                \sum_{\nPermShkVecNxt}
                                \sum_{\npermLvlVecNow}
                                \permLvlVec_{t+1}^{\koppa(\npermLvlNow,\nPermShkVecNxt)} \lomdkMat(\nmNrmVecNxt,\nmNrmVecNow,\koppa(\npermLvlNow,\nPermShkVecNxt)) \permShkRawVec(\nPermShkVecNxt) \mpMat_{t}(\nmNrmVecNow,\npermLvlVecNow) \\
                             & = 
                                \sum_{\nmNrmVecNow}
                                \sum_{\nPermShkVecNxt}
                                \permLvlVec_{t+1}^{\koppa(\npermLvlNow,\nPermShkVecNxt)} \lomdkMat(\nmNrmVecNxt,\nmNrmVecNow,\koppa(\npermLvlNow,\nPermShkVecNxt)) \permShkRawVec(\nPermShkVecNxt) \mpMatMarg_{t}(\nmNrmVecNow)
\end{aligned}\end{gathered}\end{equation}
where he second line follows because the summation of $\mpMat_{t}(\nmNrmVecNow,\npermLvlVecNow)$ is over $\npermLvlVecNow$ is exactly the step that yields $\mpMatMarg(\mNrmVecNow)$.

The steps for the permanent-income-weighted version of the proposition are identical, with the substitution of weighted for unweighted versions of the various probability objects.


\subsection{When the Joint Distribution Is Required}\label{sec:JointDistRequired}

The preceding results establish that the permanent-income-weighted distribution $\mpPrbMargWgt$ is a sufficient statistic for computing \textit{aggregate levels} of any variable $f(\mNrm)$ whose level form is $\permLvl \cdot f(\mNrm)$:
\begin{equation}\begin{gathered}\begin{aligned}
  \CLvl_{t} & = \PermGroFac^{t} \sum_{\nmNrmVecNow} \cFunc(\mNrmVecNow) \mpMatMargWgt[\nmNrmVecNow], &
  \ALvl_{t} & = \PermGroFac^{t} \sum_{\nmNrmVecNow} \aNrm(\mNrmVecNow) \mpMatMargWgt[\nmNrmVecNow]
\end{aligned}\end{gathered}\end{equation}
and analogously for aggregate income $\YLvl$ and market resources $\MLvl$. These are all \textit{linear in} $\permLvl$, which is why Harmenberg's dimensional reduction works.

However, many economically important statistics are \textit{not} linear in $\permLvl$. Computing such statistics requires the full joint distribution $\mpPrb(\mNrm, \permLvlPrb)$.

\paragraph{Variance.}
The simplest inequality measure is the cross-sectional variance of wealth levels $\ALvl_i = \permLvlPrb_i \cdot \aNrm(\mNrm_i)$:
\begin{equation}\begin{gathered}\begin{aligned} \label{eq:VarA}
  \textup{Var}(\ALvl) & = \mathbb{E}[(\permLvlPrb \cdot \aNrm)^{2}] - (\mathbb{E}[\permLvlPrb \cdot \aNrm])^{2}
\end{aligned}\end{gathered}\end{equation}
The second term can be computed from $\mpPrbMargWgt$ (it is the square of a $\permLvl$-linear aggregate). But the first term involves $\mathbb{E}[\permLvlPrb^{2} \cdot \aNrm^{2}]$, which is \textit{quadratic} in $\permLvlPrb$ and cannot be factorized via the neutral measure.

\paragraph{Quantiles and the Gini coefficient.}
Similarly, the wealth share held by the top 10\% requires sorting agents by the level $\ALvl_i = \permLvlPrb_i \cdot \aNrm_i$; the rank ordering depends jointly on $\permLvlPrb_i$ and $\mNrm_i$. The Gini coefficient:
\begin{equation}\begin{gathered}\begin{aligned} \label{eq:GiniA}
\textup{Gini}(\ALvl) & = \frac{1}{2\bar{\ALvl}} \, \mathbb{E}\!\left[|\permLvlPrb_i \aNrm_i - \permLvlPrb_j \aNrm_j|\right]
\end{aligned}\end{gathered}\end{equation}
involves the absolute difference of two agents' level wealth, which is a non-factorizable function of both agents' $(\permLvlPrb, \mNrm)$ states.

\begin{claim}\label{claim:NonAffine}
Let $h:\mathbb{R}_{+} \to \mathbb{R}$ be any function that is not affine. Then the statistic $\mathbb{E}[h(\permLvlPrb \cdot f(\mNrm))]$ generally cannot be computed from $\mpPrbMargWgt$ and $\mathbb{E}[\permLvlPrb]$ alone.
\end{claim}
\begin{proof}
By Jensen's inequality, for non-affine $h$, $\mathbb{E}[h(\permLvlPrb \cdot f)] \neq h(\mathbb{E}[\permLvlPrb] \cdot \tilde{\mathbb{E}}[f])$ in general. The Harmenberg factorization eliminates the covariance between $\permLvlPrb$ and $f(\mNrm)$ for computing means, but does not eliminate the effect of the nonlinearity of $h$ acting on the product $\permLvlPrb \cdot f$.
\end{proof}

\paragraph{Summary of $\permLvl$-linearity boundary.}
\begin{center}
\begin{tabular}{lll}
\hline
Statistic & Depends on & $\mpPrbMargWgt$ sufficient? \\
\hline
Mean consumption $\CLvl$ & $\mathbb{E}[\permLvlPrb \cdot \cFunc]$ & \textbf{Yes} \\
Mean income $\YLvl$ & $\mathbb{E}[\permLvlPrb \cdot \tranShkAll]$ & \textbf{Yes} \\
Mean assets $\ALvl$ & $\mathbb{E}[\permLvlPrb \cdot \aNrm]$ & \textbf{Yes} \\
Multiplier $\Delta\CLvl/\Delta\YLvl$ & Ratio of linear quantities & \textbf{Yes} \\
Variance of consumption & $\mathbb{E}[\permLvlPrb^{2} \cFunc^{2}]$ & \textbf{No} \\
Gini of wealth & Sorting by $\permLvlPrb \cdot \aNrm$ & \textbf{No} \\
Top 10\% wealth share & Quantiles of $\permLvlPrb \cdot \aNrm$ & \textbf{No} \\
Welfare (CRRA utility) & $\mathbb{E}[\permLvlPrb^{1-\CRRA}\cFunc^{1-\CRRA}]$ & \textbf{No} \\
\hline
\end{tabular}
\end{center}

When the joint distribution is needed, three approaches are available:
\begin{enumerate}
\item \textbf{Standard 2D transition matrix} on the $({\mNrm}, \permLvlPrb)$ grid, with cost $O((K \cdot N_{p})^{2})$.
\item \textbf{Standard Monte Carlo}, which naturally produces the joint sample $\{(\mNrm_i, \permLvlPrb_i)\}$.
\item \textbf{Harmenberg + reconstruction}: compute aggregates via $\mpPrbMargWgt$, then reconstruct the $\permLvlPrb$-dimension from the model's known age-conditional lognormal structure.
\end{enumerate}


\subsection{The Covariance Kernel and $1$D Computability}\label{sec:CovKernel}

Although inequality statistics require the full joint distribution, the \textit{covariance} between the consumption ratio $\cNrm$ and permanent income $\permLvlPrb$ can be computed exactly from the 1D weighted distribution $\mpPrbMargWgt$. This covariance appears in the aggregate consumption decomposition (cf.\ the body text, Section~\ref{subsec:Covariances}):
\begin{equation}\begin{gathered}\begin{aligned} \label{eq:CdecompApndx}
  \CLvl_{t} & = N\left(\cNrmAvg_{t} \cdot \bar{\permLvlPrb}_{t} + \cov_{t}(\cNrm, \permLvlPrb)\right)
\end{aligned}\end{gathered}\end{equation}
where $\cNrmAvg_{t}$ is the population mean consumption ratio and $\bar{\permLvlPrb}_{t}$ is the mean permanent income level.

\paragraph{Law of total covariance.}
Conditioning on the pre-shock state $(\aNrm_{t-1}, \permLvlPrb_{t-1})$:
\begin{equation}\begin{gathered}\begin{aligned} \label{eq:LTC}
\cov(\permLvlPrb_{t}, \cNrm_{t}) & = \underbrace{\mathbb{E}\!\left[\cov(\permLvlPrb_{t}, \cNrm_{t} \mid \aNrm_{t-1}, \permLvlPrb_{t-1})\right]}_{(A)} + \underbrace{\cov\!\left(\mathbb{E}[\permLvlPrb_{t} \mid \cdot], \mathbb{E}[\cNrm_{t} \mid \cdot]\right)}_{(B) = 0}
\end{aligned}\end{gathered}\end{equation}
Term~(B) vanishes in the ergodic distribution because $\permLvlPrb_{t-1}$ and $(\aNrm_{t-1})$ are asymptotically independent in the normalized dynamics.

\paragraph{The conditional covariance: $\permLvlPrb$ factors out.}
For an agent with pre-shock state $(\aNrm, \permLvlPrb)$ who transitions to permanent income $\permLvlPrb' = \permLvlPrb \cdot \PermGroFac \cdot \permShk$, the conditional covariance over shocks $(\permShk, \tranShkAll)$ is:
\begin{equation}\begin{gathered}\begin{aligned} \label{eq:CondCov}
\cov(\permLvlPrb', \cNrm' \mid \aNrm, \permLvlPrb) & = \permLvlPrb \cdot \PermGroFac \cdot \cov_{\permShk,\tranShkAll}\!\left(\permShk, \; \cFunc\!\left(\frac{\Rfree\aNrm}{\PermGroFac\permShk} + \tranShkAll\right)\right)
\end{aligned}\end{gathered}\end{equation}
The key fact is that $\permLvlPrb$ factors out linearly. Averaging over the ergodic population and using asymptotic independence of $\permLvlPrb$ and $\aNrm$:
\begin{equation}\begin{gathered}\begin{aligned} \label{eq:CovFrom1D}
\cov(\permLvlPrb_{t}, \cNrm_{t}) & = \bar{\permLvlPrb}_{t} \cdot \PermGroFac \cdot \mathbb{E}_{\aNrm}\!\left[\gamma(\aNrm)\right]
\end{aligned}\end{gathered}\end{equation}
where $\gamma(\aNrm)$ is the \textbf{covariance kernel}:
\begin{equation}\begin{gathered}\begin{aligned} \label{eq:CovKernel}
\gamma(\aNrm) & := \cov_{\permShk,\tranShkAll}\!\left(\permShk, \; \cFunc\!\left(\frac{\Rfree\aNrm}{\PermGroFac\permShk} + \tranShkAll\right)\right)
\end{aligned}\end{gathered}\end{equation}

\paragraph{$\gamma$ is computable from 1D objects.}
The covariance kernel $\gamma(\aNrm)$ depends on:
\begin{enumerate}
\item The scalar end-of-period assets $\aNrm$
\item The shock distribution $\{(\permShk_{s}, \tranShkAll_{s}, f_{s})\}$
\item The consumption function $\cFunc(\cdot)$
\end{enumerate}
None of these depend on $\permLvlPrb$. For each grid point $\aNrm_{k}$, $\gamma(\aNrm_{k})$ is computed by direct summation over the discrete shock distribution:
\begin{equation}\begin{gathered}\begin{aligned} \label{eq:GammaDiscrete}
\gamma(\aNrm_{k}) & = \sum_{s} f_{s} \cdot \permShk_{s} \cdot \cFunc\!\left(\frac{\Rfree \aNrm_{k}}{\PermGroFac \permShk_{s}} + \tranShkAll_{s}\right) \; - \; \underbrace{\left(\sum_{s} f_{s} \permShk_{s}\right)}_{=1} \underbrace{\left(\sum_{s} f_{s} \cFunc\!\left(\frac{\Rfree \aNrm_{k}}{\PermGroFac \permShk_{s}} + \tranShkAll_{s}\right)\right)}_{\mathbb{E}[\cNrm']}
\end{aligned}\end{gathered}\end{equation}

The expectation $\mathbb{E}_{\aNrm}[\gamma(\aNrm)]$ is then computed from the 1D distribution:
\begin{equation}\begin{gathered}\begin{aligned} \label{eq:EGammaTM}
\mathbb{E}_{\aNrm}[\gamma(\aNrm)] & = \sum_{k} \mpMatMargWgt[\nmNrmVecNow_{k}] \cdot \gamma(\aNrm(\mNrmVec_{t}^{k}))
\end{aligned}\end{gathered}\end{equation}
using $\mpMatMargWgt$ (since $\gamma$ is evaluated at the ergodic population).

\paragraph{Complete formula.}
Assembling \eqref{eq:CovFrom1D}, \eqref{eq:GammaDiscrete}, and \eqref{eq:EGammaTM}:
\begin{equation}\begin{gathered}\begin{aligned} \label{eq:CovComplete}
\boxed{\cov_{t}(\cNrm, \permLvlPrb) = \bar{\permLvlPrb}_{t} \cdot \PermGroFac \cdot \sum_{k} \mpMatMargWgt[\nmNrmVecNow_{k}] \cdot \gamma(\aNrm(\mNrmVec_{t}^{k}))}
\end{aligned}\end{gathered}\end{equation}
This is exact (up to grid discretization) and computed entirely from:
\begin{enumerate}
\item $\mpMatMargWgt$: the 1D permanent-income-weighted distribution (a vector of length $K$)
\item $\cFunc(\cdot)$: the consumption function
\item $\aNrm(\cdot)$: the end-of-period asset function
\item $\{(\permShk_{s}, \tranShkAll_{s}, f_{s})\}$: the shock distribution
\item $\bar{\permLvlPrb}_{t}$: mean permanent income (available analytically)
\end{enumerate}
No 2D transition matrix or 2D distribution is needed.

\paragraph{Computational cost.}
Evaluating \eqref{eq:CovComplete} costs $O(K \cdot S)$ operations (where $K$ is the number of $\mNrm$-grid points and $S$ is the number of shock realizations). Constructing a 2D transition matrix costs $O((K \cdot N_{\permLvlPrb})^{2})$. With typical values $K = 200$, $S = 50$, $N_{\permLvlPrb} = 50$: the 1D computation costs $\sim 10{,}000$ operations versus $\sim 10^{8}$ for the 2D TM.

\paragraph{Connection to balanced growth.}
In the Szeidl-invariant economy ($\cNrmAvg$ constant), $\cov_{t+1} = \PermGroFac \cdot \cov_{t}$ (see body text equation~\eqref{eq:cNrmvsCov}). More generally, in a Harmenberg-invariant economy with $\cNrmAvg$ growth factor $\GroFac_{\cNrmAvg}$, the covariance growth factor is $\GroFac_{\cov} = \GroFac_{\cNrmAvg} \cdot \PermGroFac$ (equation~\eqref{eq:aNrmGrovscovGro}). The serial correlation of $\cov_{t}$ is exactly~1 in the invariant state: knowing $\cov_{t}$ determines $\cov_{t+1}$ through a deterministic multiplicative factor. Formula~\eqref{eq:CovComplete} provides an independent way to verify this balanced-growth property numerically, using only the 1D distribution at each period.


\subsection{Higher-Order Moments from Analytical $\permLvlPrb$-Structure}\label{sec:HigherMoments}

For the inequality statistics in Section~\ref{sec:JointDistRequired}, the relevant objects involve $\mathbb{E}[\permLvlPrb^{k} \cdot g(\mNrm)]$ for $k \geq 2$. In the ergodic distribution, because $\permLvlPrb$ and $\mNrm$ are approximately independent:
\begin{equation}\begin{gathered}\begin{aligned} \label{eq:pkApprox}
\mathbb{E}[\permLvlPrb^{k} \cdot g(\mNrm)] & \approx \mathbb{E}[\permLvlPrb^{k}] \cdot \mathbb{E}[g(\mNrm)]
\end{aligned}\end{gathered}\end{equation}
where the approximation error is $O(\cov(\permLvlPrb^{k-1}, g(\mNrm)))$, which for $k=1$ is exactly the covariance computed in Section~\ref{sec:CovKernel}.

The moments $\mathbb{E}[\permLvlPrb^{k}]$ can be computed analytically from the model's permanent income dynamics. Under Blanchard--Yaari perpetual youth with survival probability $\LivPrb$ and permanent growth factor $\PermGroFac$, the second moment satisfies:
\begin{equation}\begin{gathered}\begin{aligned} \label{eq:Ep2}
\mathbb{E}[\permLvlPrb^{2}] & = \mathbb{E}[\permLvlPrb_{\textup{init}}^{2}] \cdot \frac{1 - \LivPrb}{1 - \LivPrb \cdot g_{2}}
\end{aligned}\end{gathered}\end{equation}
where $g_{2} = \PermGroFac^{2} \cdot \mathbb{E}[\permShk^{2}]$ is the per-period growth factor for $\mathbb{E}[\permLvlPrb^{2}]$.

This suggests an approximation strategy for $\textup{Var}(\ALvl)$: use 1D objects ($\mpPrbMargWgt$, $\cFunc$) for the $\mNrm$-dimension and analytical $\permLvlPrb$-moments for the permanent income dimension. However, the approximation \eqref{eq:pkApprox} is not exact for $k \geq 2$ because the mixed terms $\cov(\permLvlPrb^{k-1}, g(\mNrm))$ are not absorbed by the neutral measure.


\ifSubfilesClassLoaded{\bibfilesfind{\texname}\bibliography{\bibfilesfound}}{}

\end{document}



\endinput

\begin{comment}
One way to think about Harmenberg's key insight is to note that if we rewrite \eqref{eq:MeasTrnst} as an explicit function of the transitory and permanent shocks then the integrand becomes
\begin{equation}\begin{gathered}\begin{aligned} 
%  \mpPrb_{t+1}(\mNrmNxt,\permShkPrbNxt\pLvl_{t}) & =
                                           \lomdkPrb_{t}\left(\overbrace{(\mNrmNow-\cNrm_{t})/(\Rfree \permShkPrbNxt) + \permShkPrbNxt\tranShkNxt}^{\mNrmNxt},\mNrmNow,
\permShkPrbNxt\right)\permShkRawPrb(\permShkPrbNxt)\mpPrb_{t}(\mNrmNow,\permLvlPrbNow)/(\PermGroFac \permShkPrbNxt) \label{eq:integrand}
\end{aligned}\end{gathered}\end{equation}
or 
\begin{equation}\begin{gathered}\begin{aligned} 
\tilde{\lomdkPrb}_{t}\left(\tranShkNxt(\mNrmNxt),\mNrmNow,
\permShkPrbNxt\right)\permShkRawPrb(\permShkPrbNxt)\mpPrb_{t}(\mNrmNow,\permLvlNow)/(\PermGroFac \permShkPrbNxt) \label{eq:integrand}
\end{aligned}\end{gathered}\end{equation}

But since neither of the first two terms in the integrand \eqref{eq:integrand} involves $\pLvl$ we can change the order of integration 
\begin{equation}\begin{gathered}\begin{aligned} 
\mpPrbMarg_{t+1}(\mNrmNxt) & = \int_{\mNrmNow} \left[
  \int_{\permShkPrbNxt}  \left(
  \int_{\tranShkNxt} 
  \tilde{\lomdkPrb}_{t}\left(\tranShkNxt,\mNrmNow,
  \permShkPrbNxt\right)
  \mpPrbMarg_{t}(\mNrmNow) \Prb_{\tranShkNxt}(\tranShkNxt)d \tranShkNxt  
  \right) \permShkRawPrb(\permShkPrbNxt) d \permShkPrbNxt\right]
\mathsf{}                                                        d \mNrmNow
\end{aligned}\end{gathered}\end{equation}

\begin{equation}\begin{gathered}\begin{aligned}
 \log \GPFacRaw & < \sum_{\npermLvlVecNow} \sum_{\npermLvlVecNxt} \log\left( \permShkPrbNxt[\nPermShkNxtFunc] \right) \permShkPrbNxt[\nPermShkNxtFunc] \Prb_{\permShk}(\permShkPrbNxt[\nPermShkNxtFunc]) 
\end{aligned}\end{gathered}\end{equation}

\end{comment}

\pagebreak
Definition 2:  The permanent-income-weighted distribution is defined as 
\newcommand{\mpPrbMargAlt}{\tilde{\mpPrb}^{\mNrm}}
\begin{equation}\begin{gathered}\begin{aligned}
  \mpPrbMargAlt_{t}(\mNrmNow) & = \PermGroFac^{-t} \int_{\permLvl} \permLvlNow \mpPrb_{t}(\mNrmNow,\permLvlNow) d\permLvlNow
\end{aligned}\end{gathered}\end{equation}

\begin{equation}\begin{gathered}\begin{aligned}
  \mpPrbMargAlt_{t}[\nmNrmVecNow] & = \PermGroFac^{-t} \sum_{\npermLvlVecNow} \permLvlNow[\npermLvlVecNow] \mpPrb_{t}[\nmNrmVecNow,\npermLvlVecNow] 
\end{aligned}\end{gathered}\end{equation}

So aggregate consumption is
\begin{equation}\begin{gathered}\begin{aligned}
  \CLvl_{t} & = \int_{\permLvlNow} \left(\int_{\mNrmNow} \cNrm(\mNrmNow) \permLvlNow \mpPrb_{t}(\mNrmNow,\permLvlNow)  d  \mNrmNow\right) d \permLvlNow \\
  l\CLvl_{t} & = \int_{\permLvlNow} \left(\int_{\mNrmNow} \log (\cNrm(\mNrmNow) \permLvlNow) \mpPrb_{t}(\mNrmNow,\permLvlNow)  d  \mNrmNow\right) d \permLvlNow \\
  \\  l\CLvl_{t} & = \int_{\mNrmNow} \underbrace{\left(\int_{\permLvlNow} (\log \cNrm(\mNrmNow)+\log \permLvlNow) \mpPrb_{t}(\mNrmNow,\permLvlNow) d \permLvlNow   \right)}_{\equiv \ell \mpPrbMargAlt(\mNrmNow)} d\mNrmNow
  \\  l\CLvl_{t} & = \int_{\mNrmNow} \left(\log \cNrm(\mNrmNow)+\underbrace{\left(\int_{\permLvlNow} (\log \permLvlNow)    \right)}_{\equiv \ell \mpPrbMargAlt(\mNrmNow)} \right) \mpPrb_{t}(\mNrmNow,\permLvlNow) d \permLvlNow d\mNrmNow
  \\             & = \int_{\mNrmNow} \left(\int_{\permLvlNow} \cNrm(\mNrmNow) \permLvlNow \mpPrb_{t}(\mNrmNow,\permLvlNow) d \permLvlNow   \right) d\mNrmNow           \\   
\\            & = \int_{\mNrmNow} \cNrm(\mNrmNow)\underbrace{\left(\int_{\permLvlNow} \permLvlNow \mpPrb_{t}(\mNrmNow,\permLvlNow) d \permLvlNow   \right)}_{\equiv \mpPrbMargAlt(\mNrmNow)\PermGroFac^{t}} d\mNrmNow
\\  \CLvl_{t} & = \PermGroFac^{t} \int_{\mNrmNow} \cNrm_{t} (\mNrmNow) \mpPrbMargAlt(\mNrmNow) d\mNrmNow
\end{aligned}\end{gathered}\end{equation}
or 
\begin{equation}\begin{gathered}\begin{aligned}
  \CLvl_{t}  & = \PermGroFac^{t}\sum_{\nmNrmVecNow} \cNrm(\mNrmNow[\nmNrmVecNow])\mpPrbMargAlt[\nmNrmVecNow]
\end{aligned}\end{gathered}\end{equation}

Theorem 1:  The laws of motion for $\mpPrbMarg$ and $\mpPrbMargAlt$ are given by 
\begin{equation}\begin{gathered}\begin{aligned} 
  \mpPrbMarg_{t+1}(\mNrmNxt) & = \int_{\permShkPrbNxt} \left[
                               \int_{\mNrmNow} \lomdkPrb_{t}
                               \left(
                               \mNrmNxt,\mNrmNow,\permShkPrbNxt\right)
                               \permShkRawPrb(\permShkPrbNxt)\mpPrbMarg_{t}(\mNrmNow)
                               d \mNrmNow
                               \right]
                               d \permShkPrbNxt
\end{aligned}\end{gathered}\end{equation}

\begin{equation}\begin{gathered}\begin{aligned} 
  \mpPrbMargAlt_{t+1}(\mNrmNxt) & = \int_{\permShkPrbNxt} \left[
                               \int_{\mNrmNow} \lomdkPrb_{t}
                               \left(
                               \mNrmNxt,\mNrmNow,\permShkPrbNxt\right)
                               \permShkWgtPrb(\permShkPrbNxt)\mpPrbMargAlt_{t}(\mNrmNow)
                               d \mNrmNow
                               \right]
                               d \permShkPrbNxt
\end{aligned}\end{gathered}\end{equation}

Proposition 2:  There is a stable invariant permanent-income-weighted distribution $\mpPrbMargAlt$ if
\begin{equation}\begin{gathered}\begin{aligned}
  \log[\Rfree (1- \MPC)] & < \tilde{\Ex}\log[\PermGroFac \permShkPrbNxt]
\\ \log \GPFacRaw & < \tilde{\Ex}\log[ \permShkPrbNxt]
\\ \log \GPFacRaw & < \int_{\permShk} \log[ \permShkPrbNxt] \permShkPrbNxt f(\permShkPrbNxt) d \permShkPrbNxt
\end{aligned}\end{gathered}\end{equation}


\pagebreak
\begin{equation}\begin{gathered}\begin{aligned}
  l\bar{\mNrm}_{t} & = \PermGroFac^{t} \int_{\mNrm} l\bar{\mNrm} \mpPrbMargAlt(\mNrm) d \mNrm
  \\ & =                    \PermGroFac^{t} \tilde{\mathbb{M}}[l\mNrm]
\end{aligned}\end{gathered}\end{equation}

\begin{equation}\begin{gathered}\begin{aligned}
  & = \int_{\permLvlNow} \left(\int_{\mNrmNow} \log \left(\aNrm(\mNrmNow) \permLvlNow\right) \mpPrb_{t}(\mNrmNow,\permLvlNow)  d  \mNrmNow\right) d \permLvlNow \\
 & = \int_{\mNrmNow} \left(\int_{\permLvlNow} \log\left(\aNrm(\mNrmNow) \permLvlNow\right) \mpPrb_{t}(\mNrmNow,\permLvlNow) d \permLvlNow   \right) d\mNrmNow           \\   
 & = \int_{\mNrmNow} \log \aNrm(\mNrmNow)\underbrace{\left(\int_{\permLvlNow} (\log \permLvlNow) \mpPrb_{t}(\mNrmNow,\permLvlNow) d \permLvlNow   \right)}_{\equiv \ell \mpPrbMargAlt(\mNrmNow)\log \PermGroFac^{t}} d\mNrmNow              
\end{aligned}\end{gathered}\end{equation}

\begin{equation}\begin{gathered}\begin{aligned}
  \tilde{\Ex}[(\log \aNrm_{t+1} - \log \PermGroFac \aNrm_{t})] &
 = \Ex\left[ \tilde{\Ex}[(\log \aNrm_{t+1} - \log \PermGroFac \aNrm_{t})] + \right]
\end{aligned}\end{gathered}\end{equation}




\begin{equation}\begin{gathered}\begin{aligned} 
  \mpPrb_{t+1}[\nmNrmVecNxt,\npermLvl] & = \sum_{\nPermShk} \left[\sum_{\nmNrmVecNow} \lomdkMat_{t}\left(\mNrmNxt,\mNrmNow,\overbrace{\frac{\permLvlNxt}{\PermGroFac \permLvlNow}}^{\permShkPrbNxt}\right)\permShkRawPrb(\frac{\permLvlNxt}{\PermGroFac \permLvlNow})\mpPrb_{t}(\mNrmNow,\permLvlNow)  d \mNrmNow \right] d \permShkPrbNxt  
\end{aligned}\end{gathered}\end{equation}


% Local Variables:
% eval: (setq TeX-command-list  (remove '("Biber" "biber %s" TeX-run-Biber nil  (plain-tex-mode latex-mode doctex-mode ams-tex-mode texinfo-mode)  :help "Run Biber") TeX-command-list))
% eval: (setq TeX-command-list  (remove '("Biber" "biber %s" TeX-run-Biber nil  t  :help "Run Biber") TeX-command-list))
% eval: (setq TeX-command-list  (remove '("BibTeX" "%(bibtex) %s"    TeX-run-BibTeX nil t :help "Run BibTeX") TeX-command-list))
% eval: (setq TeX-command-list  (remove '("BibTeX" "bibtex %s"    TeX-run-BibTeX nil t :help "Run BibTeX") TeX-command-list))
% tex-bibtex-command: "bibtex.*"
% TeX-PDF-mode: t
% TeX-file-line-error: t
% TeX-debug-warnings: t
% LaTeX-command-style: (("" "%(PDF)%(latex) %(file-line-error) %(extraopts) -output-directory=. %S%(PDFout)"))
% TeX-source-correlate-mode: t
% TeX-parse-self: t
% eval: (cond ((string-equal system-type "darwin") (progn (setq TeX-view-program-list '(("Skim" "/Applications/Skim.app/Contents/SharedSupport/displayline -b %n %o %b"))))))
% eval: (cond ((string-equal system-type "gnu/linux") (progn (setq TeX-view-program-list '(("Evince" "evince --page-index=%(outpage).%o"))))))
% eval: (cond ((string-equal system-type "gnu/linux") (progn (setq TeX-view-program-selection '((output-pdf "Evince"))))))
% TeX-parse-all-errors: t
% End:
